% ================== ЛАБОРАТОРНАЯ РАБОТА №1 ==================
\documentclass[14pt,a4paper]{extarticle}

\usepackage[utf8]{inputenc}
\usepackage[russian]{babel}
\usepackage{geometry}
\geometry{left=30mm,right=15mm,top=20mm,bottom=20mm}

\usepackage{indentfirst}
\setlength{\parindent}{1.25cm}

\usepackage{amsmath,amssymb}
\usepackage{graphicx}
\usepackage{float}
\usepackage{xcolor}
\usepackage{hyperref}
\usepackage{caption}
\usepackage{setspace}
\setstretch{1.0}

\usepackage{titlesec}
\titleformat{\section}{\normalfont\normalsize\bfseries}{\thesection}{1em}{\MakeUppercase}
\titleformat{\subsection}{\normalfont\normalsize\bfseries}{\thesubsection}{1em}{}
\titlespacing{\section}{\parindent}{3.5ex plus 1ex minus .2ex}{2.3ex plus .2ex}
\titlespacing{\subsection}{\parindent}{3.25ex plus 1ex minus .2ex}{1.5ex plus .2ex}

\usepackage{fancyhdr}
\pagestyle{fancy}
\fancyhf{}
\fancyfoot[R]{\thepage}
\renewcommand{\headrulewidth}{0pt}
\fancypagestyle{plain}{
  \fancyhf{}
  \fancyfoot[R]{\thepage}
  \renewcommand{\headrulewidth}{0pt}
}

\renewcommand{\contentsname}{\centering СОДЕРЖАНИЕ}

\begin{document}

\thispagestyle{empty}

\begin{center}
МИНИСТЕРСТВО ОБРАЗОВАНИЯ РЕСПУБЛИКИ БЕЛАРУСЬ\\[0.5em]
Учреждение образования\\[0.5em]
\textbf{<<Белорусский государственный университет информатики и радиоэлектроники>>}\\[2em]

Факультет компьютерных систем и сетей\\[0.5em]
Кафедра программного обеспечения информационных технологий
\end{center}

\vspace{3em}

\begin{center}
\textbf{ОТЧЁТ}\\[0.5em]
по лабораторной работе №1\\[0.5em]
по дисциплине\\[0.5em]
\textbf{<<Генетические и эволюционные вычисления>>}

\textbf{Тема: <<Классификация с помощью персептрона>>}
\end{center}

\vspace{4em}

\hfill\begin{minipage}{0.55\textwidth}
Выполнил:\\
магистрант 1 курса магистратуры\\
факультета компьютерных систем и сетей\\
группы 556301\\[0.3em]
\textbf{Лебедевич Артём Владимирович}\\[1.2em]
Преподаватель:\\
доцент кафедры ПОИТ, канд. техн. наук\\
\textbf{Нестеренков Сергей Николаевич}
\end{minipage}

\vfill

\begin{center}
Минск, 2026~г.
\end{center}

\newpage

\tableofcontents

\newpage

\section{Цель работы}

Изучить принципы работы персептрона для решения задач классификации. Реализовать однослойный персептрон с одним и двумя нейронами, продемонстрировать его возможности и ограничения на задачах бинарной и многоклассовой классификации.

\section{Теоретические сведения}

\subsection{Персептрон}

Персептрон~--- это простейшая модель искусственного нейрона, предложенная Ф.~Розенблаттом в 1958~г. Персептрон с $n$ входами вычисляет взвешенную сумму входных сигналов и применяет к ней пороговую функцию активации:

\begin{equation}
y = f\left(\sum_{i=1}^{n} w_i x_i + b\right),
\end{equation}

где $w_i$~--- веса, $b$~--- смещение (bias), а $f$~--- функция активации (ступенчатая функция Хевисайда):

\begin{equation}
f(z) = \begin{cases} 1, & z \geq 0, \\ 0, & z < 0. \end{cases}
\end{equation}

\subsection{Алгоритм обучения персептрона}

Обучение персептрона происходит итеративно по правилу:
\begin{equation}
w_i^{(t+1)} = w_i^{(t)} + \eta \cdot (d - y) \cdot x_i,
\end{equation}
\begin{equation}
b^{(t+1)} = b^{(t)} + \eta \cdot (d - y),
\end{equation}

где $\eta$~--- скорость обучения, $d$~--- целевое значение, $y$~--- выход персептрона.

Согласно теореме о сходимости персептрона, если данные линейно разделимы, алгоритм гарантированно сходится за конечное число итераций.

\subsection{Проблема XOR}

Персептрон способен разделять только линейно разделимые классы. Функция XOR (исключающее ИЛИ) не является линейно разделимой, что было показано М.~Минским и С.~Пейпертом в 1969~г. Это ограничение стало одной из причин <<зимы>> нейронных сетей.

\section{Ход работы}

\subsection{Задание 1: Персептрон для классификации на 2 класса}

Реализован персептрон с 2 входами и 1 нейроном для классификации следующих входных векторов:

\begin{equation}
P = \begin{pmatrix} -0.5 & -0.5 & 0.3 & -0.1 \\ -0.5 & 0.5 & -0.5 & 1.0 \end{pmatrix}, \quad
T = \begin{pmatrix} 0 & 0 & 1 & 1 \end{pmatrix}.
\end{equation}

Реализован алгоритм обучения персептрона. Веса инициализировались нулями, скорость обучения $\eta = 1.0$. Обучение прошло успешно~--- персептрон сошёлся за 3 эпохи, достигнув 100\% точности.

Результаты визуализированы: входные точки окрашены по классам, построена разделяющая прямая $w_1 x_1 + w_2 x_2 + b = 0$.

\subsection{Задание 2: Проблема XOR}

Для демонстрации проблемы XOR использованы входные векторы:

\begin{equation}
P = \begin{pmatrix} 0 & 0 & 1 & 1 \\ 0 & 1 & 0 & 1 \end{pmatrix}, \quad
T = \begin{pmatrix} 0 & 1 & 1 & 0 \end{pmatrix}.
\end{equation}

Однослойный персептрон не смог сойтись за 50 эпох: ошибка осциллировала и не достигала нуля. Это подтверждает теоретический результат~--- XOR не является линейно разделимой задачей, и ни одна прямая на плоскости не способна разделить точки на заданные классы.

\subsection{Задание 3: Персептрон из двух нейронов для 4 классов}

Реализован персептрон с 2 входами и 2 выходными нейронами, каждый из которых вычисляет свою разделяющую границу. Комбинации выходов двух нейронов $(y_1, y_2)$ кодируют 4 класса: $(0,0)$, $(0,1)$, $(1,0)$, $(1,1)$.

Сгенерированы 4 кластера точек на плоскости. Персептрон обучен и сошёлся за 2 эпохи, достигнув 100\% точности. Визуализация показывает две разделяющие прямые, делящие плоскость на 4 области, каждая из которых соответствует одному классу.

\section{Результаты}

\begin{enumerate}
\item Персептрон с одним нейроном успешно решает задачу бинарной классификации линейно разделимых данных, сходясь за малое число эпох.
\item Однослойный персептрон не способен решить задачу XOR~--- ошибка не уменьшается, что подтверждает ограничения линейного классификатора.
\item Персептрон с двумя нейронами расширяет возможности классификации до 4 классов, формируя 2 разделяющие гиперплоскости.
\end{enumerate}

\section{Выводы}

В ходе лабораторной работы изучены принципы работы и обучения персептрона. Экспериментальные результаты подтвердили теоретические положения: теорему о сходимости персептрона для линейно разделимых данных и невозможность решения задачи XOR однослойной архитектурой. Многонейронный персептрон позволяет решать задачи с большим числом классов, формируя несколько разделяющих границ.

\end{document}
